================================================================================
RIGOROUS MATHEMATICAL VERIFICATION FRAMEWORK
Constrained Emergent Geometric Field via Lookup Table
================================================================================

Notation:
  - All indices are purely mathematical (no semantic labels).
  - Bodies are indexed i, j, k, l ∈ {1, ..., 6}.
  - Geometric structures are indexed by κ ∈ {1, 2, 3, 4}.
  - Tensor summation convention is active for repeated spatial indices.


--------------------------------------------------------------------------------
1. STATE SPACE AND CONFIGURATION
--------------------------------------------------------------------------------

Definition V1 (State Space).
The configuration manifold is Σ = R³⁶, coordinatized by the state vector

    q = (r₁, r₂, r₃, r₄, r₅, r₆, p₁, p₂, p₃, p₄, p₅, p₆)  ∈ Σ,

where for each i ∈ {1, ..., 6}:

    rᵢ = (rᵢ¹, rᵢ², rᵢ³) ∈ R³      (position of body i)
    pᵢ = (pᵢ¹, pᵢ², pᵢ³) ∈ R³      (canonical momentum of body i)

The symplectic form on Σ is

    ω = Σᵢ₌₁⁶  dpᵢᵃ ∧ drᵢᵃ                        (E1)

with the standard symplectic matrix J ∈ R³⁶ˣ³⁶ satisfying J² = −I.


Definition V2 (Separation Vector and Norm).
For any pair (i, j) with i < j, define the separation vector

    rᵢⱼ := rᵢ − rⱼ ∈ R³                            (E2)

and its Euclidean norm

    |rᵢⱼ| := √(δₐᵦ rᵢⱼᵃ rᵢⱼᵇ)                     (E3)

where δₐᵦ is the Kronecker delta on R³.


--------------------------------------------------------------------------------
2. THE HAMILTONIAN SYSTEM
--------------------------------------------------------------------------------

Definition V3 (Hamiltonian on Σ).
The Hamiltonian function H : Σ → R is defined as

    H(q) = T(p) + U⁽²⁾(r) + U⁽³⁾(r) + U⁽≥⁴⁾(r, p)    (E4)

where the constituent terms are:

(a) Kinetic energy:

    T(p) = Σᵢ₌₁⁶  (pᵢ · pᵢ) / (2mᵢ)               (E5)

(b) Pairwise potential energy:

    U⁽²⁾(r) = − Σ₁≤ᵢ<ⱼ≤₆  G mᵢ mⱼ / |rᵢⱼ|        (E6)

(c) Three-body correction:

    U⁽³⁾(r) = Σ₁≤ᵢ<ⱼ<ₖ≤₆  Qᵢⱼₖ / (|rᵢⱼ|² |rⱼₖ|²)  (E7)

(d) Higher-order terms from the emergent geometric field:

    U⁽≥⁴⁾(r, p) = Σₗ₌₄⁶  Uₗ(r, p)                  (E8)

with each Uₗ collecting all l-body interactions and velocity-dependent
contributions from the compactified fiber.

Here mᵢ > 0 are mass parameters, G > 0 is the coupling constant, and Qᵢⱼₖ
are three-body coupling coefficients.


Definition V4 (Hamilton's Equations).
The dynamics on Σ are governed by

    ṙᵢ =  ∂H/∂pᵢ  = pᵢ/mᵢ + ∂U⁽≥⁴⁾/∂pᵢ          (E9)
    ṗᵢ = −∂H/∂rᵢ  = −∂U⁽²⁾/∂rᵢ − ∂U⁽³⁾/∂rᵢ − ∂U⁽≥⁴⁾/∂rᵢ   (E10)

for each i ∈ {1, ..., 6}.


--------------------------------------------------------------------------------
3. THE FLOW MAP (LOOKUP TABLE)
--------------------------------------------------------------------------------

Definition V5 (Time-t Flow Map).
The time-t flow map of the Hamiltonian vector field X_H = (ṙ, ṗ) is

    Φ_Hᵗ : Σ → Σ,     Φ_Hᵗ(q₀) = q(t)             (E11)

where q(t) is the unique solution to Hamilton's equations (E9)–(E10) with
initial condition q(0) = q₀. By Liouville's theorem,

    Φ_Hᵗ preserves the symplectic form:  (Φ_Hᵗ)*ω = ω.   (E12)

The map Φ_Hᵗ is the fundamental lookup table: for any initial state q₀,
it returns the evolved state at time t.


--------------------------------------------------------------------------------
4. ERROR FUNCTIONAL AND CONSTRAINT EQUATIONS
--------------------------------------------------------------------------------

Definition V6 (Observed Trajectory Data).
Let q_obs : [0, T] → Σ denote observed trajectory data, where for each
t ∈ [0, T],

    q_obs(t) = (r₁^{obs}(t), ..., r₆^{obs}(t), p₁^{obs}(t), ..., p₆^{obs}(t)).


Definition V7 (L² Error Functional).
The error functional E[Φ_H] : {flow maps} → R≥₀ is

    E[Φ_H] = || Φ_Hᵗ(q₀) − q_obs(t) ||_{L²[0,T]}     (E13)

    = [ ∫₀ᵀ || Φ_Hᵗ(q₀) − q_obs(t) ||²_Σ dt ]^{1/2}   (E14)

where the norm on Σ is

    ||q||²_Σ = Σᵢ₌₁⁶ ( mᵢ δₐᵦ rᵢᵃ rᵢᵇ + δₐᵦ pᵢᵃ pᵢᵇ / mᵢ ).   (E15)


Proposition 1 (Euler-Lagrange Equations for E[Φ]).
The flow map Φ_H minimizes the error functional E[Φ] if and only if the
following first-order stationarity conditions hold for all t ∈ [0, T]:

    δE/δH = 0.                                     (E16)

Explicitly, let q(t) = Φ_Hᵗ(q₀) and define the residual

    η(t) := q(t) − q_obs(t) ∈ Σ.                   (E17)

Then (E16) is equivalent to

    ∫₀ᵀ ⟨ η(t),  δX_H(q(t)) ⟩_Σ dt = 0            (E18)

for all admissible variations δX_H of the Hamiltonian vector field, where
⟨·,·⟩_Σ is the inner product inducing the norm (E15).


Definition V8 (Constraint Equations for the Emergent Field).
The emergent geometric field must satisfy the following constraint system
at every t ∈ [0, T]:

(C1)   rᵢ(t) − rᵢ^{obs}(t) = 0,        ∀ i ∈ {1, ..., 6}
(C2)   pᵢ(t) − pᵢ^{obs}(t) = 0,        ∀ i ∈ {1, ..., 6}
(C3)   Σᵢ mᵢ rᵢ(t) = 0                 (center-of-mass constraint)
(C4)   Σᵢ pᵢ(t) = 0                    (momentum constraint)

In integrated form, these imply

    E[Φ_H]² = ∫₀ᵀ Σᵢ [ mᵢ |rᵢ(t) − rᵢ^{obs}(t)|² + |pᵢ(t) − pᵢ^{obs}(t)|²/mᵢ ] dt = 0.   (E19)


Corollary 1 (Determination of Coupling Constants).
The coupling parameters (mᵢ, G, Qᵢⱼₖ, and higher-order coefficients) are
constrained by:

    ∂E[Φ_H]/∂G = 0,                                (E20)
    ∂E[Φ_H]/∂Qᵢⱼₖ = 0,      ∀ (i,j,k),            (E21)
    ∂E[Φ_H]/∂mᵢ = 0,        ∀ i ∈ {1,...,6}.       (E22)

Together, (E20)–(E22) yield a closed nonlinear system for the parameters.


--------------------------------------------------------------------------------
5. CONTRIBUTIONS OF THE FOUR EMERGENT STRUCTURES
--------------------------------------------------------------------------------

The 4-dimensional emergent submanifold M with geometric structures
κ = 1, 2, 3, 4 contributes to the Hamiltonian as follows.


Proposition 2 (Structure κ = 1: Modified Pairwise Potential).
The first emergent structure modifies the pairwise potential via a geometric
correction factor. The effective two-body potential becomes

    U_{κ=1}⁽²⁾(r) = − Σᵢ<ⱼ  G mᵢ mⱼ / |rᵢⱼ|  ·  [ 1 + α₁/|rᵢⱼ| + α₂/|rᵢⱼ|² + ... ]    (E23)

where α₁, α₂, ... are dimensionless coupling constants arising from the
κ = 1 fiber geometry. In closed form:

    U_{κ=1}⁽²⁾(r) = − Σᵢ<ⱼ  (G mᵢ mⱼ / |rᵢⱼ|) · f_{κ=1}(|rᵢⱼ|/L₁)    (E24)

where L₁ is a length scale from the compactified fiber and f_{κ=1} is a
smooth dimensionless function with f_{κ=1}(0) = 1.


Proposition 3 (Structure κ = 2: Velocity-Dependent Terms).
The second emergent structure introduces velocity-dependent interactions,
breaking the strict separation T(p) + U(r). The correction is

    U_{κ=2}(r, p) = Σᵢ<ⱼ  (β₁/|rᵢⱼ|) (pᵢ · pⱼ)/(mᵢ mⱼ c²)
                  + Σᵢ<ⱼ  (β₂/|rᵢⱼ|²) [(rᵢⱼ · pᵢ)(rᵢⱼ · pⱼ)]/(mᵢ mⱼ c²)
                  + O(c⁻⁴)                                          (E25)

where β₁, β₂ are dimensionless parameters from the κ = 2 fiber geometry
coupling, and c is a fundamental speed parameter. These terms arise from
Lorentz-structure corrections in the compactified geometry.


Proposition 4 (Structure κ = 3: Three-Body Corrections).
The third emergent structure contributes to the three-body interaction term:

    U_{κ=3}⁽³⁾(r) = Σᵢ<ⱼ<ₖ  Qᵢⱼₖ^{(κ=3)} / (|rᵢⱼ|² |rⱼₖ|²)          (E26)

where the coupling coefficient decomposes as

    Qᵢⱼₖ^{(κ=3)} = γ₁ mᵢ mⱼ mₖ + γ₂ (mᵢ + mⱼ + mₖ) + γ₃.          (E27)

The constants γ₁, γ₂, γ₃ encode the κ = 3 fiber curvature contributions.
This term vanishes when any two bodies coincide, regularizing the collision
set in the configuration space.


Proposition 5 (Structure κ = 4: Higher-Order Corrections).
The fourth emergent structure generates four-body and higher interactions:

    U_{κ=4}^{(≥4)}(r) = Σᵢ<ⱼ<ₖ<ₗ  Wᵢⱼₖₗ / (|rᵢⱼ|² |rⱼₖ|² |rₖₗ|²)
                        + Σᵢ<ⱼ<ₖ<ₗ<m  Vᵢⱼₖₗₘ / (|rᵢⱼ|² |rⱼₖ|² |rₖₗ|² |rₗₘ|²)
                        + O(6-body)                                (E28)

where Wᵢⱼₖₗ and Vᵢⱼₖₗₘ are coupling tensors determined by the κ = 4
fiber topology. These terms encode the entanglement of multiple bodies
through the emergent geometric field.


The full Hamiltonian incorporating all four structures is therefore:

    H_{full}(q) = Σᵢ pᵢ²/(2mᵢ)
                + U_{κ=1}⁽²⁾(r)
                + U_{κ=2}(r, p)
                + [U_{κ=3}⁽³⁾(r) + original U⁽³⁾(r)]
                + U_{κ=4}^{(≥4)}(r)
                + U_{fiber}(r, p)                                   (E29)

where U_{fiber}(r, p) collects residual fiber-curvature effects.


--------------------------------------------------------------------------------
6. VERIFICATION CONDITION
--------------------------------------------------------------------------------

Definition V9 (Verification Condition).
Let ε > 0 be a specified error bound. The emergent geometric field is said
to verify the lookup table Φ_H to accuracy ε if and only if

    E[Φ_H] < ε.                                    (E30)

Equivalently, in terms of the pointwise residual:

    sup_{t∈[0,T]} ||Φ_Hᵗ(q₀) − q_obs(t)||_Σ < ε/√T.   (E31)

Definition V10 (Convergent Verification).
A sequence of Hamiltonians {Hₙ} with corresponding flow maps {Φ_{Hₙ}ᵗ}
converges to the observed lookup table if

    lim_{n→∞} E[Φ_{Hₙ}] = 0.                       (E32)

This is achieved when the coupling parameters of all four emergent
structures are determined by (E20)–(E22).


--------------------------------------------------------------------------------
7. EFFECTIVE GEOMETRIC POTENTIAL EQUATION
--------------------------------------------------------------------------------

Definition V11 (Mass Density Field).
The mass density field on R³ induced by the body configurations is

    ρ(r, t) = Σᵢ₌₁⁶ mᵢ δ³(r − rᵢ(t)).              (E33)


Definition V12 (Effective Geometric Potential).
The effective geometric potential Φ_eff : R³ × [0,T] → R satisfies the
nonlinear wave equation:

    ∇² Φ_eff = 4πG ρ + (1/c²) ∂ₜ² Φ_eff + Λ_eff.  (E34)

Here:

    ∇² = δᵃᵇ ∂ₐ ∂ᵦ   is the spatial Laplacian on R³,
    ∂ₜ² = ∂²/∂t²     is the second time derivative,
    Λ_eff            is the effective curvature term from the fiber.


Proposition 6 (Fiber Curvature Term Λ_eff).
The effective curvature term Λ_eff decomposes over the four emergent
structures as

    Λ_eff = Λ_{κ=1} + Λ_{κ=2} + Λ_{κ=3} + Λ_{κ=4}   (E35)

where:

(a) κ = 1 contribution (pairwise geometric modification):

    Λ_{κ=1}(r, t) = − (1/2) Σᵢ<ⱼ  G mᵢ mⱼ f''_{κ=1}(|r − rᵢⱼ̄|/L₁) / L₁²   (E36)

(b) κ = 2 contribution (velocity-dependent source):

    Λ_{κ=2}(r, t) = (4πG/c²) Σᵢ mᵢ |ṙᵢ(t)|² δ³(r − rᵢ(t))      (E37)

(c) κ = 3 contribution (three-body curvature):

    Λ_{κ=3}(r, t) = Σᵢ<ⱼ<ₖ  Qᵢⱼₖ^{(κ=3)}  K₃(r; rᵢ, rⱼ, rₖ)   (E38)

where K₃ is the kernel:

    K₃(r; rᵢ, rⱼ, rₖ) = −4 ∇² [ 1/(|rᵢⱼ|² |rⱼₖ|²) ] · δ³(r − r̄ᵢⱼₖ)

and r̄ᵢⱼₖ = (rᵢ + rⱼ + rₖ)/3.

(d) κ = 4 contribution (higher-order topology):

    Λ_{κ=4}(r, t) = Σ_{n≥4}  (−1)ⁿ λₙ  R^{(n)}(r; {rᵢ}_{i=1}⁶)   (E39)

where R^{(n)} denotes the n-th order Riemann curvature invariant of the
emergent submanifold M evaluated at the body positions, and λₙ are
normalization constants from dimensional reduction.


Theorem 1 (Self-Consistency of the Verification Framework).
Let Φ_eff be the solution to (E34) with the decomposition (E35)–(E39).
Let H_{full} be the Hamiltonian (E29). Then:

    E[Φ_{H_{full}}] < ε                                    (E40)

if and only if the following coupled system has a solution:

    (i)   Hamilton's equations for H_{full} yield q(t) = Φ_{H_{full}}ᵗ(q₀),
    (ii)  Φ_eff satisfies (E34) with source ρ from (E33),
    (iii) The coupling parameters satisfy (E20)–(E22),
    (iv)  The verification bound (E30) holds.


Proof Sketch.
(⇒) If E[Φ_{H_{full}}] < ε, then by definition (E13), the trajectory
q(t) = Φ_{H_{full}}ᵗ(q₀) is ε-close to q_obs(t) in L². The mass density
(E33) generates Φ_eff via (E34), and by the coupling constraints
(E20)–(E22), the fiber curvature terms Λ_{κ} self-consistently reproduce
the higher-order structure of H_{full}.

(⇐) Conversely, given a solution to (E34) with the decomposed Λ_eff,
the reconstructed Hamiltonian H_{full} generates a flow map. Conditions
(iii) ensure this flow map minimizes the error functional, and (iv)
guarantees the bound. ∎


--------------------------------------------------------------------------------
8. SUMMARY OF THE CONSTRAINT SYSTEM
--------------------------------------------------------------------------------

The complete constraint system for the emergent geometric field is:

    ┌─────────────────────────────────────────────────────────────────────┐
    │                        CONSTRAINT SYSTEM                            │
    ├─────────────────────────────────────────────────────────────────────┤
    │  C1.  State constraints:     rᵢ(t) = rᵢ^{obs}(t),  pᵢ(t) = pᵢ^{obs}(t)    │
    │                                                                     │
    │  C2.  Conservation:          Σᵢ mᵢ rᵢ(t) = 0,    Σᵢ pᵢ(t) = 0      │
    │                                                                     │
    │  C3.  Coupling equations:    ∂E/∂G = 0, ∂E/∂Qᵢⱼₖ = 0, ∂E/∂mᵢ = 0 │
    │                                                                     │
    │  C4.  Field equation:        ∇²Φ_eff = 4πGρ + c⁻² ∂ₜ²Φ_eff + Λ_eff│
    │                                                                     │
    │  C5.  Curvature term:        Λ_eff = Σ_{κ=1}⁴ Λ_{κ}               │
    │                                                                     │
    │  C6.  Verification bound:    E[Φ_H] < ε                           │
    │                                                                     │
    │  C7.  Structure contributions:                                      │
    │       κ=1: modified pairwise potential U_{κ=1}⁽²⁾                  │
    │       κ=2: velocity-dependent terms U_{κ=2}(r,p)                   │
    │       κ=3: three-body correction U_{κ=3}⁽³⁾                        │
    │       κ=4: higher-order interaction U_{κ=4}^{(≥4)}                 │
    └─────────────────────────────────────────────────────────────────────┘

The four emergent geometric structures (κ = 1, 2, 3, 4) contribute to
the Hamiltonian in a hierarchical manner:

    - κ = 1 modifies the fundamental two-body interaction, encoding
      fiber-scale geometric corrections to the inverse-distance law.

    - κ = 2 introduces relativistic velocity-dependent couplings,
      reflecting the Lorentz structure of the compactified geometry.

    - κ = 3 regularizes the dynamics via three-body interactions that
      dominate at short separation distances.

    - κ = 4 generates the full topological complexity of the emergent
      submanifold through multi-body entanglement.

The effective geometric potential Φ_eff self-consistently ties all
contributions together through the wave equation (E34), with the fiber
curvature Λ_eff acting as a distributed source encoding the compactified
dimension data. The verification condition E[Φ_H] < ε closes the system,
ensuring that the emergent field reproduces the lookup table Φ_Hᵗ to
the prescribed accuracy.

================================================================================
                              END OF FRAMEWORK
================================================================================
